\section{Discussion}\label{sec:discussion}

\subsection{Principal Findings}

This systematic literature review examined 120 studies of computational methods for angiographic sequence analysis. In this review, we found that publications increased rapidly after 2020, with 78\% appearing from 2022 onward. However, shared evaluation, reproducibility, and clinical validation remain limited in the reviewed corpus. The main findings of this manuscript have been summarized in four key takeaways:

First, temporal information is commonly used, but the way it is modeled varies substantially. Because eligibility required temporally ordered multi-frame data, these proportions characterize the selected sequence-analysis literature rather than angiographic AI as a whole. Although 97 studies (80.8\%) used sequential data, 49 use fixed operations and only 48 learn temporal representations. Learned temporal models are concentrated in classification tasks and remain scarce in flow analysis, device tracking, and outcome prediction. Moreover, 41 studies use short frame stacks rather than complete acquisitions, and 51 of the 97 studies that used sequential data still identify limited temporal modeling as a limitation. As shown in Table~\ref{tab:temporal_taxonomy}, using sequence data therefore covers approaches with very different temporal capabilities. The next research step is therefore to evaluate when learned temporal representations improve on simpler methods through controlled comparisons using the same data, inputs, and evaluation protocol.

Second, the literature is concentrated in a small number of clinical settings. Cardiac and neurologic applications constitute nearly 90\% of the included studies, and only five studies use both DSA and fluoroscopy. Evidence from these settings does not establish performance in other vascular territories or acquisition conditions. Broader evaluation should therefore include differences in anatomy, imaging vendors, frame rates, and acquisition protocols, rather than relying only on additional data from the development setting.

Third, reported performance is generally high, but its interpretation is limited by inconsistent evaluation. Public datasets are used by 20 studies, yet only three release their own data, and there is no common evaluation framework across the reviewed tasks. Differences in datasets, reference standards, and metrics make cross-study rankings unreliable. Therefore, public data benchmarks need to be built with defined test sets and consistent outcome definitions.

Fourth, evidence supporting clinical use remains limited. Only 19 studies report external validation, three use a prospective design, and four report clinical deployment or live use. These findings indicate that most evidence concerns technical performance rather than benefit during clinical care. Similarly, the reporting comparison shown in Figure~\ref{fig:forest_plot} found no statistically significant differences between studies published up to 2024 and those published from 2025 onward after Holm correction. This does not establish unchanged reporting or measure the effect of the CLAIM 2024 update.

The twelve research gaps bring these findings together. Table~\ref{tab:research_gaps} shows that model development is closely connected to data availability, annotation quality, evaluation, and clinical integration. Generalization, public benchmarks, and temporal modeling are the most frequently documented gaps, appearing in 68, 65, and 63 studies, respectively. These provide immediate priorities for further work. Federated learning, language models, and structured reporting are less represented and should be treated as directions requiring evaluation rather than established solutions.

\subsection{Comparison with Prior Work}

The surveys by Litjens et al.~\cite{Litjens2017} and Shen et al.~\cite{Shen2017} examine deep learning across medical imaging, while Moccia et al.~\cite{Moccia2018} review blood vessel segmentation, datasets, and evaluation metrics. Our review complements these perspectives by focusing on temporally ordered X-ray angiographic data and distinguishing fixed temporal operations from learned temporal representations. Using multiple frames may involve aggregation, registration, flow-curve fitting, or learning dependencies across a sequence; these approaches provide different temporal capabilities. Their frequency in the reviewed literature does not establish their effectiveness. Whether learned temporal models improve performance or clinical usefulness requires controlled comparisons using the same data, inputs, and evaluation protocol.

The reviews by Iezzi et al.~\cite{Iezzi2019}, Gurgitano et al.~\cite{Gurgitano2021}, and Mukhtar et al.~\cite{Mukhtar2025} examine AI across interventional radiology, including planning, guidance, interpretation, and clinical workflow. Our contribution is to examine data availability, external validation, evaluation, and clinical implementation within a defined corpus of angiographic sequence studies and relate them to tasks, architectures, and temporal modeling. We also assess reporting using review-specific criteria informed by CLAIM 2024~\cite{Tejani2024} (Appendix~\ref{app:claim_clusters}). Background and objectives were more consistently reported than demographics and subgroup performance; separately, code and data release were uncommon. Meeting the broader reproducibility and funding-disclosure cluster (C12) does not necessarily indicate artifact release. These findings distinguish reporting completeness and reproducibility resources from methodological validity and evidence of clinical usefulness.

\subsection{Limitations}
\label{sec:limitations}

First, the search covered only PubMed/\allowbreak MEDLINE and Semantic Scholar, up to 25 May 2026. PubMed searches were restricted to English, and each Semantic Scholar query was limited to its first 100 results. Conference papers and arXiv-only preprints may therefore be underrepresented. Publication counts for 2026 also cover only part of the year and should not be compared directly with full-year counts.

Second, title and abstract screening relied on an LLM. Two independent reviewers checked all 134 pipeline inclusions and a random sample of 357 exclusions, but most excluded records were not reviewed manually. Two eligible studies were found in the exclusion sample, giving an estimated 28 eligible records among the 4,930 original pipeline exclusions (95\% Wilson interval scaled to that stratum: 8--100), including the two records recovered by human review. Both failed the sequence-data criterion because their abstracts did not explicitly establish eligibility. The exclusion-by-default rules for I2 and I4 may therefore have reduced recall. The corpus is a validated sample of the retrieved literature rather than a complete census, and missed studies could affect category proportions as well as counts.

Third, full-text extraction relied on an LLM, and one reviewer audited six studies, representing 5\% of the corpus. Although 97.1\% of audited values were correct, this sample cannot establish the same accuracy for every field. Errors in interpretive categories, including temporal modeling and research gaps, may affect the synthesis. The gap analysis also depends on limitations and future directions stated by authors, so problems that were not explicitly discussed may be underrepresented.

Fourth, the reporting assessment used an operational schema informed by CLAIM 2024. As documented in Appendix~\ref{app:claim_clusters}, its 44 criteria do not consistently correspond to the numbering and definitions of the official checklist. The 12 thematic groups were defined for this review, and a group was counted as met when at least half of its applicable items were reported. These results should therefore be interpreted as reporting completeness under our schema, rather than strict adherence to the official CLAIM checklist. Group-level summaries may also conceal missing individual items. With 43 studies in the later publication period, the absence of statistically significant differences does not establish that reporting remained unchanged or that the guideline had no impact.

Fifth, we did not apply a dedicated risk-of-bias tool or conduct a meta-analysis. The reported performance distributions describe the extracted results but do not provide pooled estimates of generalization or clinical effectiveness. Differences in study settings and evaluation methods also prevent conclusions about whether one architecture or temporal modeling strategy is superior.

Finally, the review was not registered and had no published protocol. Full text was unavailable for two eligible studies, and three included records are earlier versions of another included study (R2434/R1238, R4154/R3555, and R4671/R4479 in \texttt{extraction/manifest.csv}). These related records introduce dependence and slightly inflate some counts.

\subsection{Conclusion}

This review shows growing interest in angiographic sequence analysis, with applications spanning classification, detection, segmentation, tracking, and image enhancement. Temporal information is widely used, but learned temporal representations remain concentrated in a few tasks. High reported performance is accompanied by limited data sharing, inconsistent evaluation, and little prospective evidence of clinical benefit.

Further progress requires shared multi-institutional datasets, task-specific benchmarks, controlled evaluation of temporal methods, and validation across vascular territories and acquisition settings. Prospective studies should then establish whether these methods improve clinical decisions, procedural workflow, or patient outcomes. These steps are necessary to turn technical improvements into reliable tools for angiography-guided care.
